% Sample Applications, Chapter 2 -- Loan Payment Table (annex Sec 8.3).
% Values replayed through the shipped application (examples/loan.i23);
% @PMT as os/i123/core/eval.c computes it (graded against the real
% Release 2.2 formulas by test-i123-formula).  Figure screens/sa2.k123.
% GATES: test-i123-formula test-pubs-i123-examples test-pubs-screens
\Gates{test-i123-formula, test-pubs-i123-examples, test-pubs-screens}
\chapter{Loan Payment Table}

Before the purchasing office commits to financing new equipment, it
compares the monthly payment on the loan at several interest rates.  This
worksheet computes the payment for a given principal and term at four
annual rates.

\section{What You Enter}
\begin{itemize}
\item The principal, the amount borrowed
\item The term of the loan in years
\item The annual interest rates to compare
\end{itemize}

\section{Entering the Data}
\ix{"@PMT}
Type the labels \type{Principal} in A3 and \type{Years} in A4, and the
amounts beside them in B3 and B4.  Type the headings \type{"Rate} and
\type{"Payment} in A6 and B6, and the rates in A7..A10 as decimals,
\type{.08} for 8 percent.  In B7 enter the payment formula
\type{@pmt(\$b\$3,a7/12,\$b\$4*12)} and copy it to B8..B10.  Format the rates
as percentages and the payments with two decimal places.
\FlairShot[fig:sa2]{sa2}{Monthly payments at four rates}

\begin{LotTable}{p{30mm}p{24mm}p{74mm}}
\LotHead{\textbf{Item} & \textbf{Address} & \textbf{Meaning}}
Principal & B3 & The amount borrowed\\
Years & B4 & The term of the loan\\
Rate & A7..A10 & The annual interest rates\\
Payment & B7..B10 & The monthly payment: \texttt{@PMT} of the principal at
  the monthly rate (the annual rate divided by 12) over the term in
  months (years times 12)\LotEnd
\end{LotTable}

\section{Interpreting the Results}
The payment formulas and the payments:
\LotExample{examples/loan}
The principal and the term are absolute references, so every copy of the
formula uses them; the rate is relative, so each row uses its own.

\section{Changing the Template}
Type a new principal in B3 or a new term in B4, and every payment is
recalculated.  To compare more rates, type them below A10 and copy the
formula of B10 down beside them.
